
\documentclass[11pt,a4paper]{article}
\usepackage[utf8]{inputenc}
\usepackage[T1]{fontenc}
\usepackage[french]{babel}
\usepackage{lmodern}
\usepackage{microtype}
\usepackage[a4paper,margin=1.65cm,headheight=15pt]{geometry}
\usepackage{amsmath,amssymb,mathtools}
\usepackage{array,tabularx,multicol,booktabs}
\usepackage{enumitem}
\usepackage[most]{tcolorbox}
\usepackage{xcolor}
\usepackage{tikz}
\usetikzlibrary{arrows.meta,calc,positioning}
\usepackage{pdflscape}
\usepackage{fancyhdr}
\usepackage{hyperref}
\usepackage{needspace}
\usepackage{pagecolor}

\definecolor{fondpage}{RGB}{247,248,250}
\definecolor{encre}{RGB}{31,35,40}
\definecolor{bleu}{RGB}{40,91,135}
\definecolor{vert}{RGB}{55,125,92}
\definecolor{orange}{RGB}{178,105,42}
\definecolor{violet}{RGB}{101,77,145}
\definecolor{rouge}{RGB}{170,72,72}
\definecolor{gris}{RGB}{86,91,99}
\definecolor{bleufonce}{RGB}{28,55,120}
\definecolor{bleuclair}{RGB}{238,243,255}
\definecolor{grisclair}{RGB}{245,245,245}
\definecolor{tableline}{RGB}{45,45,45}
\definecolor{softgray}{RGB}{246,247,249}
\definecolor{deepblue}{RGB}{25,62,115}
\definecolor{accent}{RGB}{20,82,130}

\hypersetup{colorlinks=true,linkcolor=bleufonce,urlcolor=bleufonce,
pdftitle={Parcours de revision -- Fonctions trigonométriques -- Premiere specialite},pdfauthor={}}

\setlength{\parindent}{0pt}
\setlength{\parskip}{0.25em}
\renewcommand{\arraystretch}{1.13}
\setlength{\tabcolsep}{4pt}
\setlist[itemize]{leftmargin=1.25em,itemsep=0.15em,topsep=0.2em}
\setlist[enumerate,1]{label=\textbf{\arabic*.}, leftmargin=1.05cm, itemsep=0.35em, topsep=0.2em}
\setlist[enumerate,2]{label=\textbf{\alph*.}, leftmargin=0.85cm, itemsep=0.25em, topsep=0.15em}

\pagestyle{fancy}
\fancyhf{}
\cfoot{\thepage}

\newcommand{\R}{\mathbb{R}}
\newcommand{\N}{\mathbb{N}}
\newcommand{\sourceq}[1]{ {\scriptsize\textcolor{bleufonce}{(site Premiere q#1)}}}
\newcommand{\code}[1]{\texttt{#1}}

\newcounter{question}
\newcounter{reponse}

\newenvironment{blocquestions}[1]{%
  \begin{tcolorbox}[enhanced,breakable,colback=bleuclair,colframe=bleufonce,boxrule=0.6pt,arc=2mm,left=2mm,right=2mm,top=2mm,bottom=1.5mm,title={\bfseries #1},coltitle=white,colbacktitle=bleufonce,before skip=3mm,after skip=3mm, fontupper=\large]
  \large
  \begin{enumerate}[leftmargin=*,label=\textbf{\arabic*.},itemsep=0.82em,topsep=0.5em]
  \setcounter{enumi}{\value{question}}
}{%
  \setcounter{question}{\value{enumi}}
  \end{enumerate}
  \end{tcolorbox}
}

\newenvironment{blocreponses}[1]{%
  \begin{tcolorbox}[enhanced,breakable,colback=bleuclair,colframe=bleufonce,boxrule=0.6pt,arc=2mm,left=2mm,right=2mm,top=2mm,bottom=1.5mm,title={\bfseries #1},coltitle=white,colbacktitle=bleufonce,before skip=3mm,after skip=3mm, fontupper=\large]
  \large
  \begin{enumerate}[leftmargin=*,label=\textbf{\arabic*.},itemsep=0.42em,topsep=0.5em]
  \setcounter{enumi}{\value{reponse}}
}{%
  \setcounter{reponse}{\value{enumi}}
  \end{enumerate}
  \end{tcolorbox}
}

\newtcolorbox{correctionbox}[1]{enhanced,breakable,colback=grisclair,colframe=black!55,boxrule=0.55pt,arc=2mm,left=2mm,right=2mm,top=2mm,bottom=1.5mm,title=\textbf{#1},coltitle=black,colbacktitle=black!12,before skip=2mm,after skip=2mm}
\newtcolorbox{methodebox}[1]{enhanced,breakable,colback=white,colframe=bleu!70!black,boxrule=0.55pt,arc=2mm,left=2mm,right=2mm,top=2mm,bottom=1.5mm,title=\textbf{#1},coltitle=white,colbacktitle=bleu!70!black,before skip=2mm,after skip=2mm}

\newcommand{\titreSujet}[2]{%
  \subsection*{#1}
  \addcontentsline{toc}{subsection}{#1}
  \textit{Référence / inspiration : #2}\par\medskip\hrule\medskip
}
\newcommand{\qcm}[4]{%
\begin{center}\begin{tabularx}{0.96\linewidth}{@{}XX@{}}
\textbf{a.} #1 & \textbf{b.} #2\\[1mm]
\textbf{c.} #3 & \textbf{d.} #4
\end{tabularx}\end{center}}

\tcbset{enhanced,arc=2mm,outer arc=2mm,boxrule=0.42pt,left=1.15mm,right=1.15mm,top=0.75mm,bottom=0.75mm,boxsep=0.35mm,before skip=0.8mm,after skip=0.8mm,fonttitle=\bfseries\footnotesize,coltitle=encre,before upper=\raggedright\fontsize{7.65}{8.15}\selectfont}
\newtcolorbox{fichebox}[3][]{colback=white,colframe=#2!72!black,colbacktitle=#2!18,coltitle=#2!50!black,borderline west={0.9mm}{0pt}{#2!78!black},title={#3},drop fuzzy shadow=black!5,#1}
\newcommand{\tagcol}[2]{\begin{tcolorbox}[enhanced,colback=#1!11,colframe=#1!45!black,boxrule=0.42pt,arc=2mm,left=1mm,right=1mm,top=0.45mm,bottom=0.45mm,before skip=0mm,after skip=0.85mm,fontupper=\bfseries\scriptsize,colupper=#1!55!black]\centering #2\end{tcolorbox}}
\newtcolorbox{panel}[1]{colback=fondpage,colframe=fondpage,boxrule=0pt,arc=0mm,left=0mm,right=0mm,top=0mm,bottom=0mm,height=168mm,valign=top,before skip=0mm,after skip=0mm}
\newcommand{\titrepage}{\begin{tcolorbox}[colback=white,colframe=bleu!70!black,boxrule=0.65pt,arc=3mm,left=2.5mm,right=2.5mm,top=1.15mm,bottom=1.15mm,before skip=0mm,after skip=1.4mm,drop fuzzy shadow=black!10]
{\Large\bfseries Parcours de révision -- Fonctions trigonométriques}\hfill {\normalsize Première spécialité -- sans calculatrice}
\end{tcolorbox}}

% Tableaux TikZ sobres pour les corrections
\newcommand{\TabVarTrigF}{%
\begin{center}\begin{tikzpicture}[x=1cm,y=0.62cm,>=Latex]
\def\L{1.9}\def\W{13.8}\def\H{3.65}
\fill[black!8] (0,2.82) rectangle (\W,\H);\fill[black!5] (0,0) rectangle (\L,\H);
\draw[tableline,line width=0.45pt] (0,0) rectangle (\W,\H);\draw[tableline,line width=0.45pt] (\L,0)--(\L,\H);\draw[tableline,line width=0.45pt] (0,2.82)--(\W,2.82);\draw[tableline,line width=0.45pt] (0,1.82)--(\W,1.82);
\draw[tableline,line width=0.35pt] (5.2,0)--(5.2,\H);\draw[tableline,line width=0.35pt] (9.6,0)--(9.6,\H);
\node at (0.95,3.22) {$x$};\node at (0.95,2.32) {$f'(x)$};\node at (0.95,0.9) {$f$};
\node at (2.25,3.22) {$0$};\node at (5.2,3.22) {$\frac{\pi}{3}$};\node at (9.6,3.22) {$\frac{5\pi}{3}$};\node at (13.25,3.22) {$2\pi$};
\node at (3.6,2.32) {$+$};\node at (5.2,2.32) {$0$};\node at (7.4,2.32) {$-$};\node at (9.6,2.32) {$0$};\node at (11.45,2.32) {$+$};
\node at (2.25,0.35) {$0$};\node at (5.2,1.45) {$\sqrt3-\frac{\pi}{3}$};\node at (9.6,0.35) {$-\sqrt3-\frac{5\pi}{3}$};\node at (13.25,1.05) {$-2\pi$};
\draw[->,line width=0.65pt] (2.55,0.45)--(4.75,1.35);\draw[->,line width=0.65pt] (5.75,1.35)--(9.0,0.45);\draw[->,line width=0.65pt] (10.15,0.45)--(12.75,0.95);
\end{tikzpicture}\end{center}}

\newcommand{\TabVarH}{%
\begin{center}\begin{tikzpicture}[x=1cm,y=0.62cm,>=Latex]
\def\L{1.9}\def\W{12.2}\def\H{3.65}
\fill[black!8] (0,2.82) rectangle (\W,\H);\fill[black!5] (0,0) rectangle (\L,\H);
\draw[tableline,line width=0.45pt] (0,0) rectangle (\W,\H);\draw[tableline,line width=0.45pt] (\L,0)--(\L,\H);\draw[tableline,line width=0.45pt] (0,2.82)--(\W,2.82);\draw[tableline,line width=0.45pt] (0,1.82)--(\W,1.82);
\draw[tableline,line width=0.35pt] (7.0,0)--(7.0,\H);
\node at (0.95,3.22) {$t$};\node at (0.95,2.32) {$h'(t)$};\node at (0.95,0.9) {$h$};
\node at (2.35,3.22) {$0$};\node at (7.0,3.22) {$6$};\node at (11.65,3.22) {$12$};
\node at (4.6,2.32) {$-$};\node at (7.0,2.32) {$0$};\node at (9.25,2.32) {$+$};
\node at (2.35,1.45) {$5$};\node at (7.0,0.35) {$1$};\node at (11.65,1.45) {$5$};
\draw[->,line width=0.65pt] (2.75,1.35)--(6.5,0.45);\draw[->,line width=0.65pt] (7.5,0.45)--(11.2,1.35);
\end{tikzpicture}\end{center}}
\begin{document}
\begin{center}
{\LARGE\bfseries Corrigé -- Parcours de révision -- Fonctions trigonométriques}\\[1mm]
{\large Première générale -- spécialité mathématiques}\\[1mm]
{\normalsize Sans calculatrice}
\end{center}
\tableofcontents\newpage

\newcommand{\ficheRecap}{%
\begin{landscape}
\pagecolor{fondpage}
\thispagestyle{empty}
\titrepage
\begin{tabularx}{\linewidth}{@{}X@{\hspace{1.8mm}}X@{\hspace{1.8mm}}X@{\hspace{1.8mm}}X@{}}
\begin{panel}{}
\tagcol{bleu}{Cercle et radian}
\begin{fichebox}{bleu}{Cercle trigonométrique}
Cercle de centre $O$, de rayon $1$, orienté dans le sens direct.

Un réel $x$ correspond au point $M(x)$ de coordonnées :
\[
M(x)=(\cos x;\sin x).
\]
Un tour complet correspond à $2\pi$ radians.
\end{fichebox}
\begin{fichebox}{bleu}{Schéma repère}
\begin{center}
\begin{tikzpicture}[scale=1.05,>=Latex]
\draw[->] (-1.35,0)--(1.45,0) node[right] {$x$};
\draw[->] (0,-1.35)--(0,1.45) node[above] {$y$};
\draw[thick] (0,0) circle (1);
\draw[->,bleu!70!black,thick] (0.55,0) arc (0:42:0.55);
\draw[thick,rouge!80!black] (0,0)--(40:1) node[pos=0.6,above] {$1$};
\fill (40:1) circle (1.2pt) node[above right] {$M(x)$};
\draw[dashed] (40:1)--({cos(40)},0) node[below] {$\cos x$};
\draw[dashed] (40:1)--(0,{sin(40)}) node[left] {$\sin x$};
\node at (0.26,0.13) {$x$};
\node[below left] at (0,0) {$O$};
\end{tikzpicture}
\end{center}
\end{fichebox}
\begin{fichebox}{bleu}{Repères}
\[
\pi\ \text{rad}=180^\circ,\quad \frac{\pi}{2}=90^\circ,\quad \frac{\pi}{3}=60^\circ.
\]
\[
\frac{\pi}{4}=45^\circ,\quad \frac{\pi}{6}=30^\circ.
\]
\end{fichebox}
\end{panel}
&
\begin{panel}{}
\tagcol{vert}{Valeurs exactes}
\begin{fichebox}{vert}{Valeurs remarquables}
\centering
\begin{tabular}{c|ccccc}
$x$ & $0$ & $\frac{\pi}{6}$ & $\frac{\pi}{4}$ & $\frac{\pi}{3}$ & $\frac{\pi}{2}$ \\ \hline
$\cos x$ & $1$ & $\frac{\sqrt3}{2}$ & $\frac{\sqrt2}{2}$ & $\frac12$ & $0$ \\
$\sin x$ & $0$ & $\frac12$ & $\frac{\sqrt2}{2}$ & $\frac{\sqrt3}{2}$ & $1$
\end{tabular}
\end{fichebox}
\begin{fichebox}{vert}{Symétries utiles}
\[
\cos(-x)=\cos x,\qquad \sin(-x)=-\sin x.
\]
\[
\cos(\pi-x)=-\cos x,\qquad \sin(\pi-x)=\sin x.
\]
\[
\cos(\pi+x)=-\cos x,\qquad \sin(\pi+x)=-\sin x.
\]
\end{fichebox}
\begin{fichebox}{vert}{Identité fondamentale}
Pour tout réel $x$ :
\[
\cos^2 x+\sin^2 x=1.
\]
Elle permet de retrouver un sinus à partir d’un cosinus, ou inversement, avec le signe donné par le quadrant.
\end{fichebox}
\end{panel}
&
\begin{panel}{}
\tagcol{orange}{Fonctions sinus et cosinus}
\begin{fichebox}{orange}{Périodicité}
Pour tout réel $x$ :
\[
\cos(x+2\pi)=\cos x,
\qquad
\sin(x+2\pi)=\sin x.
\]
Les fonctions sinus et cosinus sont $2\pi$-périodiques.
\end{fichebox}
\begin{fichebox}{orange}{Parité}
\[
\cos(-x)=\cos x \quad \Rightarrow \quad \cos \text{ est paire.}
\]
\[
\sin(-x)=-\sin x \quad \Rightarrow \quad \sin \text{ est impaire.}
\]
\end{fichebox}
\begin{fichebox}{orange}{Dérivées de référence}
\[
(\sin x)'=\cos x,
\qquad
(\cos x)'=-\sin x.
\]
Avec une composée affine :
\[
(\sin(ax+b))'=a\cos(ax+b),
\]
\[
(\cos(ax+b))'=-a\sin(ax+b).
\]
\end{fichebox}
\end{panel}
&
\begin{panel}{}
\tagcol{violet}{Méthodes}
\begin{fichebox}{violet}{Résoudre une équation simple}
Pour résoudre $\cos x=a$ ou $\sin x=a$ sur un intervalle :
\begin{itemize}
\item placer la valeur sur le cercle ;
\item repérer les angles associés ;
\item respecter l’intervalle demandé ;
\item ajouter $2k\pi$ si on résout dans $\R$.
\end{itemize}
\end{fichebox}
\begin{fichebox}{violet}{Étudier une fonction trigonométrique}
\begin{enumerate}[leftmargin=*,itemsep=0.15em,topsep=0.15em]
\item Calculer la dérivée.
\item Résoudre $f'(x)=0$.
\item Étudier le signe de $f'(x)$ sur l’intervalle.
\item Dresser le tableau de variations.
\end{enumerate}
\end{fichebox}
\begin{fichebox}{violet}{Attention}
Ne pas confondre :
\[
\cos x=\frac12 \quad \text{et}\quad x=\frac12.
\]
Le cosinus et le sinus sont des coordonnées sur le cercle trigonométrique.
\end{fichebox}
\end{panel}
\end{tabularx}
\clearpage
\nopagecolor
\end{landscape}}

\ficheRecap
\section*{Correction des questions flash}
\addcontentsline{toc}{section}{Correction des questions flash}
\setcounter{reponse}{0}
\begin{blocreponses}{Réponses 1 à 12}
\item C’est le cercle de centre $O$ et de rayon $1$, orienté dans le sens direct.
\item C’est le point de coordonnées $(\cos x;\sin x)$.
\item Un tour complet correspond à $2\pi$ radians.
\item $\cos(0)=1$.
\item $\sin(0)=0$.
\item $\cos(\pi)=-1$ et $\sin(\pi)=0$.
\item $\cos\left(\dfrac{\pi}{2}\right)=0$ et $\sin\left(\dfrac{\pi}{2}\right)=1$.
\item $\cos\left(\dfrac{\pi}{3}\right)=\dfrac12$ et $\sin\left(\dfrac{\pi}{3}\right)=\dfrac{\sqrt3}{2}$.
\item Les deux valent $\dfrac{\sqrt2}{2}$.
\item $\cos\left(\dfrac{\pi}{6}\right)=\dfrac{\sqrt3}{2}$ et $\sin\left(\dfrac{\pi}{6}\right)=\dfrac12$.
\item Pour tout réel $x$, $\cos(-x)=\cos x$.
\item Pour tout réel $x$, $\sin(-x)=-\sin x$.
\end{blocreponses}
\begin{blocreponses}{Réponses 13 à 24}
\item Pour tout réel $x$, $\cos(x+2\pi)=\cos x$ et $\sin(x+2\pi)=\sin x$.
\item On a toujours $\cos^2 x+\sin^2 x=1$.
\item $x=\dfrac{\pi}{3}$ ou $x=\dfrac{5\pi}{3}$.
\item $x=\dfrac{2\pi}{3}$ ou $x=\dfrac{4\pi}{3}$.
\item $x=\dfrac{\pi}{3}$ ou $x=\dfrac{2\pi}{3}$.
\item $x=\dfrac{7\pi}{6}$ ou $x=\dfrac{11\pi}{6}$.
\item $\cos x=-\dfrac12$, donc $x=\dfrac{2\pi}{3}$ ou $x=\dfrac{4\pi}{3}$.
\item $\sin x=\dfrac12$, donc $x=\dfrac{\pi}{6}$ ou $x=\dfrac{5\pi}{6}$.
\item $(\sin x)'=\cos x$.
\item $(\cos x)'=-\sin x$.
\item $f'(x)=3\cos x-2$.
\item $g'(x)=-2\sin x+2x$.
\end{blocreponses}
\begin{blocreponses}{Réponses 25 à 36}
\item $h'(x)=2\cos(2x)$.
\item $k'(x)=-3\sin(3x-1)$.
\item Elle est impaire, car $\sin x$ et $x$ sont impaires.
\item Non, à cause du terme $-x$.
\item $\cos(x+2\pi)=\cos x$.
\item Comme $\sin(-x)=-\sin x$, la somme vaut $0$.
\item $\dfrac{13\pi}{4}\equiv -\dfrac{3\pi}{4}\pmod{2\pi}$.
\item $x=\dfrac{\pi}{2}+2k\pi$, avec $k\in\mathbb Z$.
\item $x=\dfrac{\pi}{2}$ ou $x=\dfrac{3\pi}{2}$.
\item $\cos(t+4\pi)+\cos(-t)=\cos t+\cos t=\dfrac43$.
\item $\sin\left(\dfrac{5\pi}{6}\right)=\dfrac12$.
\item $\cos\left(\dfrac{7\pi}{6}\right)=-\dfrac{\sqrt3}{2}$.
\end{blocreponses}
\section*{Correction du QCM}
\addcontentsline{toc}{section}{Correction du QCM}
\begin{correctionbox}{Question 1}
Réponse \textbf{a.} $\dfrac12$.
\end{correctionbox}
\begin{correctionbox}{Question 2}
Réponse \textbf{a.} $\dfrac12$.
\end{correctionbox}
\begin{correctionbox}{Question 3}
Réponse \textbf{b.} paire.
\end{correctionbox}
\begin{correctionbox}{Question 4}
Réponse \textbf{b.} impaire.
\end{correctionbox}
\begin{correctionbox}{Question 5}
Réponse \textbf{b.} $-\sin x$.
\end{correctionbox}
\begin{correctionbox}{Question 6}
Réponse \textbf{b.} $\dfrac{2\pi}{3}$ et $\dfrac{4\pi}{3}$.
\end{correctionbox}
\begin{correctionbox}{Question 7}
Réponse \textbf{b.} $2\cos(2x)$.
\end{correctionbox}
\begin{correctionbox}{Question 8}
Réponse \textbf{b.} $2\pi$.
\end{correctionbox}

\section*{Correction des sujets type bac / E3C}
\addcontentsline{toc}{section}{Correction des sujets type bac / E3C}

\titreSujet{Sujet 1 -- Cercle trigonométrique et valeurs remarquables}{Correction}
\begin{correctionbox}{1. Placement des angles}
On place les angles dans le sens direct sur le cercle trigonométrique : $0$ sur l’axe des abscisses positives, puis $\frac{\pi}{6}$, $\frac{\pi}{4}$, $\frac{\pi}{3}$, $\frac{\pi}{2}$ et enfin $\pi$ sur l’axe des abscisses négatives.
\end{correctionbox}
\begin{correctionbox}{2. Valeurs exactes}
\[
\cos\left(\frac{\pi}{3}\right)=\frac12,
\qquad
\sin\left(\frac{\pi}{3}\right)=\frac{\sqrt3}{2}.
\]
Comme $\frac{5\pi}{6}=\pi-\frac{\pi}{6}$, on a :
\[
\cos\left(\frac{5\pi}{6}\right)=-\frac{\sqrt3}{2},
\qquad
\sin\left(\frac{5\pi}{6}\right)=\frac12.
\]
\end{correctionbox}
\begin{correctionbox}{3. Symétries}
Comme le cosinus est pair et le sinus est impair :
\[
\cos\left(-\frac{\pi}{4}\right)=\frac{\sqrt2}{2},
\qquad
\sin\left(-\frac{\pi}{4}\right)=-\frac{\sqrt2}{2}.
\]
De plus $\frac{7\pi}{6}=\pi+\frac{\pi}{6}$, donc :
\[
\cos\left(\frac{7\pi}{6}\right)=-\frac{\sqrt3}{2},
\qquad
\sin\left(\frac{7\pi}{6}\right)=-\frac12.
\]
\end{correctionbox}
\begin{correctionbox}{4. Équation $\cos x=\frac12$}
Sur $[0;2\pi[$, les solutions sont :
\[
x=\frac{\pi}{3}\quad \text{ou}\quad x=\frac{5\pi}{3}.
\]
\end{correctionbox}
\begin{correctionbox}{5. Équation $\sin x=-\frac{\sqrt3}{2}$}
Sur $[0;2\pi[$, les solutions sont :
\[
x=\frac{4\pi}{3}\quad \text{ou}\quad x=\frac{5\pi}{3}.
\]
\end{correctionbox}
\begin{correctionbox}{6. Périodicité}
La fonction cosinus est $2\pi$-périodique : pour tout réel $x$, $\cos(x+2\pi)=\cos x$. Les solutions se répètent donc à chaque tour complet du cercle.
\end{correctionbox}

\titreSujet{Sujet 2 -- Étude d’une fonction trigonométrique}{Correction}
\begin{correctionbox}{1. Dérivée}
Pour tout $x\in[0;2\pi]$ :
\[
f'(x)=2\cos x-1.
\]
\end{correctionbox}
\begin{correctionbox}{2. Résolution de $2\cos x-1=0$}
\[
2\cos x-1=0 \Longleftrightarrow \cos x=\frac12.
\]
Sur $[0;2\pi]$, les solutions sont :
\[
x=\frac{\pi}{3}\quad \text{ou}\quad x=\frac{5\pi}{3}.
\]
\end{correctionbox}
\begin{correctionbox}{3 à 5. Signe, variations et valeurs}
On sait que $f'(x)=2\cos x-1$. Ainsi $f'(x)$ est positif lorsque $\cos x\geq \frac12$.

\TabVarTrigF
Les valeurs utiles sont :
\[
f(0)=0,
\quad
f\left(\frac{\pi}{3}\right)=\sqrt3-\frac{\pi}{3},
\]
\[
f\left(\frac{5\pi}{3}\right)=-\sqrt3-\frac{5\pi}{3},
\quad
f(2\pi)=-2\pi.
\]
\end{correctionbox}
\begin{correctionbox}{6. Extremums locaux}
Le maximum local est atteint pour $x=\dfrac{\pi}{3}$ et vaut :
\[
\sqrt3-\frac{\pi}{3}.
\]
Le minimum local est atteint pour $x=\dfrac{5\pi}{3}$ et vaut :
\[
-\sqrt3-\frac{5\pi}{3}.
\]
\end{correctionbox}

\titreSujet{Sujet 3 -- Fonction trigonométrique en contexte périodique}{Correction}
\begin{correctionbox}{1. Valeurs exactes}
\[
h(0)=3+2\cos(0)=5.
\]
\[
h(3)=3+2\cos\left(\frac{\pi}{2}\right)=3.
\]
\[
h(6)=3+2\cos(\pi)=1.
\]
\[
h(9)=3+2\cos\left(\frac{3\pi}{2}\right)=3.
\]
\[
h(12)=3+2\cos(2\pi)=5.
\]
\end{correctionbox}
\begin{correctionbox}{2 et 3. Maximum et minimum}
Comme $-1\leq \cos\left(\frac{\pi}{6}t\right)\leq 1$, on a :
\[
1\leq h(t)\leq 5.
\]
La hauteur minimale est donc $1$ et la hauteur maximale est $5$.
\end{correctionbox}
\begin{correctionbox}{4. Période}
La fonction $t\mapsto \cos\left(\frac{\pi}{6}t\right)$ a pour période $T$ telle que :
\[
\frac{\pi}{6}T=2\pi.
\]
Donc $T=12$. La période de $h$ est $12$ secondes.
\end{correctionbox}
\begin{correctionbox}{5 à 7. Dérivée, signe et variations}
\[
h'(t)=2\times\left(-\frac{\pi}{6}\sin\left(\frac{\pi}{6}t\right)\right)
=-\frac{\pi}{3}\sin\left(\frac{\pi}{6}t\right).
\]
Sur $[0;12]$, l’angle $\frac{\pi}{6}t$ parcourt $[0;2\pi]$. Le sinus est positif sur $]0;\pi[$, donc pour $t\in]0;6[$, puis négatif sur $]\pi;2\pi[$, donc pour $t\in]6;12[$.

\TabVarH
\end{correctionbox}

\end{document}
